\documentclass[12pt,a4paper]{ctexart} % 中文技术报告模板
\usepackage[margin=2.2cm]{geometry} % 设置页边距
\usepackage{amsmath} % 数学公式支持
\usepackage{graphicx} % 图片支持
\usepackage{hyperref} % 交叉引用与链接
\usepackage[backend=biber,style=numeric]{biblatex} % 参考文献管理
\IfFileExists{examples/references.bib}{%
  \addbibresource{examples/references.bib} % 从项目根目录编译时使用这个路径
}{%
  \addbibresource{references.bib} % 从 examples/ 目录编译时使用同目录镜像
}

\title{用 LaTeX 记录一次简单实验} % 报告标题
\author{你的名字} % 作者
\date{\today} % 日期

\begin{document}

\maketitle % 生成标题区
\tableofcontents % 自动生成目录

\section{摘要} % 第 1 节
本文演示如何用 LaTeX 组织一份短小的技术报告，包括公式、图、表、交叉引用和参考文献。

\section{问题背景} % 第 2 节
本报告以一个简单练习为例，目标是熟悉 LaTeX 的基本结构，并观察“改代码”和“改版面”的对应关系。

\section{方法} % 第 3 节
我们用一个简单线性关系举例：
\begin{equation}
  y = kx + b % 这里可以换成你自己的公式
\end{equation}

\section{结果} % 第 4 节

\subsection{图示结果}
\begin{figure}[htbp]
  \centering

  % 先用占位框，保证模板不依赖真实图片也能编译
  \fbox{\rule{0pt}{4cm}\rule{0.7\linewidth}{0pt}}

  % 如果你已经准备好图片，可以改成下面这一行
  % \includegraphics[width=0.7\linewidth]{assets/figure-01.png}

  \caption{实验结果示意图} % 图题
  \label{fig:result} % 图标签
\end{figure}

图~\ref{fig:result} 展示了结果示意。 % 正文中引用图编号

\subsection{表格结果}
\begin{table}[htbp]
  \centering
  \begin{tabular}{lcc}
    \hline
    指标 & 方案A & 方案B \\
    \hline
    时间成本 & 10分钟 & 8分钟 \\
    错误次数 & 3次 & 1次 \\
    满意度 & 7分 & 9分 \\
    \hline
  \end{tabular}
  \caption{结果对比表} % 表题
  \label{tab:result} % 表标签
\end{table}

表~\ref{tab:result} 给出了一个简单对比。 % 正文中引用表编号

\section{讨论} % 第 5 节
如果需要引用文献，可以写成文献~\cite{lamport1994latex} 或文献~\cite{knuth1984texbook}。

\section{结论} % 第 6 节
通过这份模板，你已经把标题、章节、公式、图、表、交叉引用和参考文献放到了同一个文档里。

\printbibliography % 输出参考文献列表

\end{document}
